\documentclass[aps,prd,twocolumn,superscriptaddress,nofootinbib]{revtex4-2}

\usepackage{amsmath,amssymb}
\usepackage{graphicx}
\usepackage{hyperref}
\usepackage{xcolor}
\usepackage{bm}
\usepackage{multirow}
\usepackage{booktabs}

% Shortcuts (matching Papers I--II)
\newcommand{\LCDM}{$\Lambda$CDM}
\newcommand{\Lam}{\Lambda}
\newcommand{\Om}{\Omega_{\rm m}}
\newcommand{\Ob}{\Omega_{\rm b}}
\newcommand{\Or}{\Omega_{\rm r}}
\newcommand{\OL}{\Omega_\Lambda}
\newcommand{\Odom}{\Omega_{\rm dom}}
\newcommand{\fdom}{f_{\rm dom}}
\newcommand{\zc}{z_{\rm c}}
\newcommand{\sln}{\sigma_{\ln}}
\newcommand{\Geff}{G_{\rm eff}}
\newcommand{\Hubble}{H_0}
\newcommand{\dAIC}{\Delta{\rm AIC}}
\newcommand{\dchi}{\Delta\chi^2}
\newcommand{\gearly}{g_{\rm early}}
\newcommand{\glate}{g_{\rm late}}
\newcommand{\etaz}{\eta_0}
\newcommand{\Rhat}{\hat{R}}

\begin{document}

\title{Inflationary Domain Theory III: Phase Resonance, Gravitational Dephasing, and the Conserved Net Impact Observable}

\author{Jose Ramon Gonzalez}
\affiliation{Independent Researcher}

\date{\today}

% =====================================================================
% ABSTRACT
% =====================================================================
\begin{abstract}
In Papers~I and~II~\cite{Paper1,Paper2}, we introduced Inflationary
Domain Theory (IDT) as a minimal extension of $\Lambda$CDM in which
epoch-localized energy density contributions modify the expansion
history and structure growth at specific cosmological epochs, achieving
$\dAIC = -8.1$ with hand-chosen domain widths and perturbative
dephasing ($\etaz = -0.05$). In this paper, we present three advances.
First, we introduce the MS phase quantum formalism, an effective
description in which each domain resonance accumulates a constant
cosmological phase $\Delta = \int(1+w)\,dN$, providing a
data-supported relation that determines domain widths $\sln$ from a
single parameter. Second, we perform full parameter liberation with
converged posteriors from adaptive Metropolis--Hastings sampling (four
independent chains, all Gelman--Rubin $\Rhat < 1.1$, 32{,}000
post-burn-in samples). Third, we identify the combination
$g = \fdom \times (1 + \etaz)$ as a conserved observable governing the
net dephasing effect: the posterior is bimodal in $(\fdom, \etaz)$
space but unimodal in $g$, with $\gearly = 0.035 \pm 0.016$. With six
parameters beyond $\Lambda$CDM (all free, data-determined), the model
achieves $\dAIC = -24$ under the compressed Planck likelihood with
$P(k)$ constraints (AIC penalty of~12). The $S_8$ tension is reduced
from approximately $3\sigma$ to $0.2\sigma$, and the CMB acoustic
scale $l_a$ is reproduced to $0.0\sigma$. The dephasing parameter
$\etaz = -0.57$ ($-0.70$ to $-0.40$) produces behavior analogous to a
gravitational Lagrange point, where competing influences from
neighboring domains lead to partial cancellation of the effective
gravitational coupling. The bimodal posterior represents a physical
degeneracy between neighbor configurations that produce identical net
dephasing, analogous to the emergence of $\Om h^2$ as the primary
CMB-constrained quantity rather than $\Om$ and $h$ separately.
$\Hubble = 68.1 \pm 0.5\;\mathrm{km\,s^{-1}\,Mpc^{-1}}$ is
consistent with the CMB-inferred global expansion rate; the
discrepancy with local distance-ladder measurements may reflect a
local void effect rather than requiring additional modifications to
early-universe physics.
\end{abstract}

\maketitle

% =====================================================================
% SECTION 1: INTRODUCTION
% =====================================================================
\section{Introduction}
\label{sec:intro}

\subsection{The parameter motivation problem}
\label{sec:param_motivation}

In Paper~I~\cite{Paper1}, we introduced Inflationary Domain Theory
(IDT) as a minimal two-parameter extension of $\Lambda$CDM in which
epoch-localized energy density contributions modify the expansion
history at specific cosmological epochs. The framework achieved
$\dAIC = -8.1$ relative to $\Lambda$CDM, with the improvement driven
primarily by a correction to the CMB acoustic angular scale from
$3.1\sigma$ tension to $0.3\sigma$, accompanied by a shift in the
Hubble constant from $67.4$ to
$68.7\;\mathrm{km\,s^{-1}\,Mpc^{-1}}$. In Paper~II~\cite{Paper2}, we
demonstrated that the early domain ($\zc \sim 2000$) produces a
$20$--$50\%$ enhancement in high-redshift galaxy abundance at
$z = 10$--$15$, consistent with elevated abundances reported by the
James Webb Space Telescope (JWST), and mapped the $\etaz$ parameter
space across $\etaz \in [-0.05, -0.01]$.

A natural limitation of Papers~I and~II is that several key
parameters---the domain widths $\sln^{\rm early}$ and $\sln^{\rm
late}$, as well as the dephasing coupling $\etaz$---were hand-chosen
rather than constrained by the data. The domain widths in Paper~I
($\sln^{\rm early} = 0.18$, $\sln^{\rm late} = 0.50$) were selected
on physical grounds, and the dephasing coupling $\etaz = -0.05$
represented a perturbative regime explored for computational
tractability. While these choices produced statistically meaningful
improvements over $\Lambda$CDM, the question remained: what values do
the data prefer when all parameters are freed?

This question has both statistical and physical dimensions. From a
statistical perspective, freeing additional parameters increases the
AIC penalty ($2k$) and requires that any improvement in fit quality
exceed this penalty to maintain a negative $\dAIC$. From a physical
perspective, the location of the posterior maximum in the freed
parameter space carries information about the nature of the domain
interactions. If the data prefer $\etaz$ values far from the
perturbative regime explored in Papers~I--II, this would indicate that
the dephasing mechanism operates through a qualitatively different
physical process than initially assumed.

\subsection{The MS phase quantum hypothesis}
\label{sec:phase_hypothesis}

The MS phase quantum formalism provides a framework for constraining
domain widths from a single parameter rather than specifying them by
hand. The central quantity is the cosmological phase
\begin{equation}
\Phi(N) = \int_0^N (1 + w_{\rm bg})\, dN',
\label{eq:cosmo_phase}
\end{equation}
where $N = \ln a$ is the number of e-folds and $w_{\rm bg}$ is the
background equation of state. The phase velocity $d\Phi/dN = 1 + w$
tracks the rate at which the cosmological background evolves: during
radiation domination, $d\Phi/dN = 4/3$; during matter domination,
$d\Phi/dN = 1$; and during $\Lambda$-domination, $d\Phi/dN \to 0$ as
the phase freezes.

The hypothesis is that each domain resonance window accumulates a
constant phase
\begin{equation}
\Delta = \int_{N_{\rm on}}^{N_{\rm off}} (1 + w_{\rm bg})\, dN,
\label{eq:phase_quantum}
\end{equation}
where $N_{\rm on}$ and $N_{\rm off}$ define the effective boundaries
of the resonance window (determined by the Gaussian profile at, e.g.,
the $1\sigma$ level). Given $\Delta$ and the resonance center
redshift $\zc$, the domain width $\sln$ is determined by root-finding:
earlier resonances (where the phase velocity is larger) require
narrower $\sln$ to accumulate the same phase $\Delta$. This replaces
two free parameters (one $\sln$ per domain) with a single universal
constant $\Delta$, providing a data-supported relation between
resonance epoch and resonance width.

The epistemological status of this formalism is that of an effective
description. We do not claim that the constant-phase condition derives
from an underlying action principle. Rather, it organizes the
parameter space of IDT in a manner that is empirically validated by
three independent lines of evidence: grid improvement over Paper~I
(Sec.~\ref{sec:grid}), convergence of $\Delta$ under free-parameter
MCMC (Sec.~\ref{sec:progressive}), and robustness under dataset
ablation (Sec.~\ref{sec:robustness}).

\subsection{From perturbative dephasing to partial gravitational cancellation}
\label{sec:perturbative_to_partial}

In Paper~I, the dephasing coupling was set to $\etaz = -0.05$,
representing a $5\%$ modification to the effective gravitational
coupling during domain resonance. This perturbative regime was
motivated by the expectation that domain interactions should produce
small corrections to standard gravitational dynamics.

When $\etaz$ is freed in this work, the posterior converges to
$\etaz = -0.57$ ($-0.70$ to $-0.40$ at $68\%$ confidence), indicating
that the data prefer a regime of partial gravitational cancellation
rather than perturbative leakage. This behavior is analogous to a
Lagrange point in classical gravitational systems, where competing
gravitational influences from two bodies lead to a region of reduced
net acceleration. In the IDT context, the effective gravitational
coupling during domain resonance is
\begin{equation}
\Geff(z) = 1 + \etaz \cdot \Gamma(z),
\label{eq:geff_intro}
\end{equation}
where $\Gamma(z)$ encodes the domain resonance profile. At the
resonance peak, $\Geff \approx 1 + \etaz \approx 0.43$, representing
a $57\%$ reduction in gravitational coupling. This analogy is
descriptive, not constitutive: the mathematical content is
Eq.~(\ref{eq:geff_intro}), and the Lagrange point language serves to
connect the phenomenology to familiar gravitational physics.

The perturbative regime explored in Paper~I ($\etaz = -0.05$) is not
favored within the tested parameter space when $\etaz$ is freed,
indicating that the dephasing mechanism is more significant than
initially assumed.

\subsection{Paper~III scope}
\label{sec:scope}

This paper presents five connected advances:
\begin{enumerate}
\item[(a)] The MS phase quantum formalism, which provides a
  data-supported relation determining domain widths from a single
  parameter $\Delta$ (Sec.~\ref{sec:phase_formalism}).
\item[(b)] Progressive parameter liberation with converged posteriors
  from adaptive Metropolis--Hastings sampling
  (Sec.~\ref{sec:progressive}).
\item[(c)] The identification of $g = \fdom \times (1 + \etaz)$ as
  the conserved observable governing the net dephasing effect
  (Sec.~\ref{sec:posteriors}).
\item[(d)] The Lagrange point interpretation of gravitational
  dephasing at $\etaz \approx -0.57$
  (Sec.~\ref{sec:dephasing_physics}).
\item[(e)] Reduction of the $S_8$ tension from approximately $3\sigma$
  to $0.2\sigma$, CMB acoustic scale consistency to $0.0\sigma$, and
  a discussion of $\Hubble$ in the context of the local void
  hypothesis (Sec.~\ref{sec:dephasing_physics}).
\end{enumerate}
All IDT parameters are free and data-determined. All posteriors are
converged ($\Rhat < 1.1$ across four independent chains). Nothing is
hand-picked.

The paper is organized as follows. Section~\ref{sec:phase_formalism}
develops the MS phase quantum formalism.
Section~\ref{sec:grid} validates the formalism against Paper~I
through a grid comparison. Section~\ref{sec:progressive} presents
the progressive parameter liberation. Section~\ref{sec:posteriors}
analyzes the converged posteriors and the conserved observable.
Section~\ref{sec:dephasing_physics} discusses the physics of
dephasing. Section~\ref{sec:robustness} presents robustness tests.
Section~\ref{sec:predictions} outlines predictions and future tests.
Section~\ref{sec:discussion} provides a discussion of claims and
scope. Section~\ref{sec:conclusions} summarizes the conclusions.
Appendices~\ref{app:solver}--\ref{app:sigma8} provide technical
details.

% =====================================================================
% SECTION 2: MS PHASE QUANTUM FORMALISM
% =====================================================================
\section{The MS Phase Quantum Formalism}
\label{sec:phase_formalism}

\subsection{Cosmological phase and phase velocity}
\label{sec:cosmo_phase}

The cosmological phase $\Phi(N)$ defined in Eq.~(\ref{eq:cosmo_phase})
is a monotonically increasing function of the scale factor that
encodes the cumulative expansion history of the background. Its
derivative, the phase velocity
\begin{equation}
\frac{d\Phi}{dN} = 1 + w_{\rm bg}(N),
\label{eq:phase_velocity}
\end{equation}
is directly related to the equation-of-state parameter of the total
cosmological fluid. During the radiation-dominated epoch,
$w_{\rm bg} = 1/3$ and the phase velocity is $d\Phi/dN = 4/3$. During
matter domination, $w_{\rm bg} = 0$ and $d\Phi/dN = 1$. As the
cosmological constant comes to dominate, $w_{\rm bg} \to -1$ and the
phase velocity approaches zero: the cosmological phase freezes.

The transition between these regimes is smooth, with the phase
velocity passing through characteristic values at the
radiation-matter equality ($z_{\rm eq} \approx 3400$) and the
matter-$\Lambda$ equality ($z_\Lambda \approx 0.3$). Between these
transitions, $d\Phi/dN$ decreases monotonically from $4/3$ to~0,
tracing the gradual deceleration and eventual acceleration of cosmic
expansion.

For a $\Lambda$CDM background with $\Om = 0.315$, $\Or = 9.15 \times
10^{-5}$, and $\OL = 1 - \Om - \Or$, the total phase accumulated
from recombination ($z = 1100$) to the present is approximately
$\Phi \approx 7.0$, with the majority accumulated during the
matter-dominated era.

\begin{figure}[t]
\centering
\includegraphics[width=\columnwidth]{figures/fig01_phase_velocity.pdf}
\caption{Phase velocity $d\Phi/dN = 1 + w_{\rm bg}$ as a function of
  redshift for a $\Lambda$CDM background. The shaded bands indicate
  the early ($\zc \approx 2250$) and late ($\zc \approx 1.8$) domain
  resonance windows at the best-fit parameters from the full-free
  MCMC. The phase velocity is $4/3$ during radiation domination,
  transitions through $7/6$ at matter-radiation equality, approaches
  unity during matter domination, and freezes as $\Lambda$ comes to
  dominate. The early domain resonates during the radiation-to-matter
  transition; the late domain resonates as the phase velocity
  approaches zero.}
\label{fig:phase_velocity}
\end{figure}

Figure~\ref{fig:phase_velocity} shows the phase velocity as a
function of redshift, with the early and late domain resonance windows
overlaid. A notable feature is that the early domain ($\zc \approx
2250$) resonates during the radiation-to-matter transition, where the
phase velocity is intermediate ($d\Phi/dN \approx 1.1$--$1.3$), while
the late domain ($\zc \approx 1.8$) resonates as the phase velocity
is already declining toward zero ($d\Phi/dN \approx 0.8$--$1.0$).
This asymmetry has a direct consequence for the constant-phase
condition: to accumulate the same total phase $\Delta$, the early
domain (with faster phase velocity) requires a narrower resonance
window than the late domain (with slower phase velocity).

\subsection{The constant phase quantum condition}
\label{sec:constant_phase}

The central hypothesis of the MS phase quantum formalism is that each
domain resonance accumulates the same total cosmological phase:
\begin{equation}
\Delta = \int_{N_{\rm on}}^{N_{\rm off}} (1 + w_{\rm bg})\, dN
= \int_{N_{\rm on}}^{N_{\rm off}} \frac{d\Phi}{dN}\, dN,
\label{eq:phase_quantum_def}
\end{equation}
where $N_{\rm on}$ and $N_{\rm off}$ are defined by the domain
resonance profile. For a Gaussian profile in $\ln a$ centered at
$\ln a_c$ with width $\sln$, we define the resonance boundaries at
the $\pm 1\sigma$ level:
\begin{equation}
N_{\rm on} = \ln a_c - \sln, \qquad N_{\rm off} = \ln a_c + \sln.
\label{eq:resonance_bounds}
\end{equation}

Given the cosmological background (which determines $w_{\rm bg}$ as a
function of $N$), the condition in Eq.~(\ref{eq:phase_quantum_def})
establishes a one-to-one mapping between the phase quantum $\Delta$,
the resonance center $\zc$ (equivalently $\ln a_c = -\ln(1 + \zc)$),
and the domain width $\sln$. Specifically, for fixed $\Delta$ and
$\zc$, the width $\sln$ is determined by the implicit equation
\begin{equation}
\Delta = \int_{-\sln}^{+\sln} \left[1 + w_{\rm bg}(\ln a_c +
  s)\right] ds,
\label{eq:sigma_implicit}
\end{equation}
which is solved numerically by bisection (see
Appendix~\ref{app:solver}).

The physical content of this condition is as follows. If domain
resonances are triggered by a phase-matching condition between our
domain and its MS neighbors---that is, resonance occurs when the
cosmological phase in our domain matches the phase configuration of
the neighbor---then the duration of each resonance (in e-folds) should
scale inversely with the local phase velocity. This produces
resonances that accumulate equal total phase, regardless of the
cosmological epoch at which they occur.

\begin{figure}[t]
\centering
\includegraphics[width=\columnwidth]{figures/fig02_sigma_vs_zc.pdf}
\caption{Domain width $\sln$ as a function of resonance center
  redshift $\zc$ for four values of the phase quantum: $\Delta = 0.25$
  (lowest curve), $0.50$, $0.75$, and $1.00$ (highest curve). At
  fixed $\Delta$, earlier resonances (higher $\zc$, faster phase
  velocity) require narrower widths to accumulate the same total
  phase. The vertical dashed lines indicate the best-fit resonance
  centers ($\zc^{\rm early} = 2254$, $\zc^{\rm late} = 1.82$). For the
  converged value $\Delta = 0.36$, the predicted widths are
  $\sln^{\rm early} = 0.15$ and $\sln^{\rm late} = 0.19$.}
\label{fig:sigma_vs_zc}
\end{figure}

Figure~\ref{fig:sigma_vs_zc} shows $\sln$ as a function of $\zc$ for
several values of $\Delta$. Two features are apparent. First, at fixed
$\Delta$, the width increases toward lower redshifts as the phase
velocity decreases (more e-folds are needed to accumulate the same
phase). Second, the sensitivity to $\Delta$ is approximately linear:
doubling $\Delta$ roughly doubles $\sln$, with departures from strict
linearity arising from the variation of $w_{\rm bg}$ across the
resonance window.

\subsection{Parameter reduction}
\label{sec:param_reduction}

The constant-phase condition reduces the free parameter count of IDT.
In Paper~I, the model had four specified parameters ($\zc^{\rm
early}$, $\sln^{\rm early}$, $\zc^{\rm late}$, $\sln^{\rm late}$) and
two fitted parameters ($\fdom^{\rm early}$, $\fdom^{\rm late}$), for a
total of six IDT-specific parameters (with $\etaz$ hand-chosen). In
the present work, the six free IDT parameters are:
\begin{align}
\text{Domain centers:}&\quad \zc^{\rm early},\; \zc^{\rm late},
\nonumber\\
\text{Domain amplitudes:}&\quad \fdom^{\rm early},\; \fdom^{\rm late},
\nonumber\\
\text{Phase quantum:}&\quad \Delta, \nonumber\\
\text{Dephasing coupling:}&\quad \etaz.
\label{eq:free_params}
\end{align}
The domain widths $\sln^{\rm early}$ and $\sln^{\rm late}$ are derived
quantities, determined at each MCMC step from $(\Delta, \zc^{\rm
early}, \zc^{\rm late})$ via Eq.~(\ref{eq:sigma_implicit}). Together
with the three $\Lambda$CDM parameters that are varied ($\Hubble$,
$\Om$, $\Ob h^2$), the model has nine free parameters in total, six
beyond $\Lambda$CDM, yielding an AIC penalty of $2 \times 6 = 12$.

\subsection{Computational implementation}
\label{sec:implementation}

The phase quantum solver is invoked at each step of the MCMC chain.
Given the current values of $(\Hubble, \Om, \Or, \OL)$ and the
proposed $(\Delta, \zc^{\rm early}, \zc^{\rm late})$, the solver
computes $w_{\rm bg}(N)$ from the Friedmann equation and determines
$\sln^{\rm early}$ and $\sln^{\rm late}$ by bisection. The solver
converges to machine precision in $\lesssim 20$ iterations, adding
negligible computational cost relative to the CLASS Boltzmann
integration.

An important feature of the implementation is that the derived
$\sln$ values are self-consistent with the current cosmological
parameters. As the MCMC explores different values of $\Om$ and
$\Hubble$, the phase velocity profile changes, and the derived
$\sln$ values adjust accordingly. This ensures that the constant-phase
condition is maintained across the full parameter space, not just at a
fiducial cosmology.

Details of the solver algorithm, convergence properties, and
validation tests are provided in Appendix~\ref{app:solver}.

\subsection{Epistemological status of the phase formalism}
\label{sec:epistemological}

The MS phase quantum formalism should be understood as an effective
description that organizes the parameter space of IDT. It provides a
one-parameter family of domain configurations, indexed by $\Delta$,
that can be tested against data. The formalism does not derive from a
specified action principle or fundamental symmetry, and we do not
claim that the constant-phase condition is a consequence of any
particular theory of the MS substrate.

The validity of the formalism is assessed empirically through three
independent lines of evidence:

\begin{enumerate}
\item \textit{Grid improvement over Paper~I.} All MS-predicted
  $\sln$ pairs outperform the hand-chosen values of Paper~I in a
  systematic grid comparison (Sec.~\ref{sec:grid}).
\item \textit{$\Delta$ convergence under free-parameter MCMC.} When
  $\Delta$ is freed along with all other IDT parameters, the posterior
  converges to a well-defined peak at $\Delta = 0.36$ ($0.27$--$0.49$
  at $68\%$ confidence), indicating that the data select a specific
  value rather than being indifferent (Sec.~\ref{sec:progressive}).
\item \textit{Robustness under dataset ablation.} Removing the $S_8$
  constraint from the likelihood does not significantly shift the
  preferred $\Delta$, indicating that the convergence is driven by the
  CMB acoustic scale rather than the structure-growth data alone
  (Sec.~\ref{sec:robustness}).
\end{enumerate}

These lines of evidence support the utility of the phase quantum
formalism as an organizing principle for IDT parameter space, without
implying that it constitutes a fundamental law.

% =====================================================================
% SECTION 3: GRID VALIDATION
% =====================================================================
\section{Grid Validation}
\label{sec:grid}

\subsection{Design}
\label{sec:grid_design}

To test the MS phase quantum formalism against the hand-chosen
parameters of Paper~I, we perform a systematic grid comparison.
The baseline is the Paper~I configuration: $\sln^{\rm early} = 0.18$,
$\sln^{\rm late} = 0.50$, with $\fdom^{\rm early}$ and $\fdom^{\rm
late}$ optimized and $\etaz = -0.05$ fixed. The baseline achieves
$\dAIC = -8.1$~\cite{Paper1}.

We generate four MS $\sln$ pairs from the constant-phase condition
with $\Delta = 0.40$, $0.55$, $0.75$, and $1.00$, using the Paper~I
resonance centers ($\zc^{\rm early} = 2000$, $\zc^{\rm late} = 1.0$)
and the Planck best-fit cosmology. For each $\Delta$, the phase
quantum solver produces a pair $(\sln^{\rm early}, \sln^{\rm late})$,
and we optimize $\fdom^{\rm early}$ and $\fdom^{\rm late}$ by
$\chi^2$ minimization against the same likelihood used in Paper~I.
All other parameters ($\etaz$, $\Hubble$, $\Om$, $\Ob h^2$) are held
at the Paper~I values to enable a controlled comparison.

\subsection{Results}
\label{sec:grid_results}

\begin{table}[t]
\centering
\caption{Grid validation results. Each row corresponds to a different
  phase quantum $\Delta$, with the derived domain widths and optimized
  amplitudes. All four MS configurations outperform the Paper~I
  hand-chosen baseline ($\dAIC = -8.1$). The best-performing
  configuration ($\Delta = 0.75$) achieves $\dAIC = -12.8$.}
\label{tab:grid}
\begin{ruledtabular}
\begin{tabular}{lccccc}
Config & $\Delta$ & $\sln^{\rm early}$ & $\sln^{\rm late}$ &
$\dchi$ & $\dAIC$ \\
\hline
Paper~I & --- & 0.18 & 0.50 & $-12.1$ & $-8.1$ \\
MS grid & 0.40 & 0.16 & 0.21 & $-13.8$ & $-9.8$ \\
MS grid & 0.55 & 0.22 & 0.29 & $-15.4$ & $-11.4$ \\
MS grid & 0.75 & 0.30 & 0.39 & $-16.8$ & $-12.8$ \\
MS grid & 1.00 & 0.40 & 0.52 & $-15.1$ & $-11.1$ \\
\end{tabular}
\end{ruledtabular}
\end{table}

Table~\ref{tab:grid} summarizes the grid results. All four MS
configurations outperform the Paper~I baseline, with $\dAIC$ ranging
from $-9.8$ to $-12.8$. The best-performing configuration has
$\Delta = 0.75$, achieving $\dAIC = -12.8$---an improvement of $4.7$
units over Paper~I.

\begin{figure}[t]
\centering
\includegraphics[width=\columnwidth]{figures/fig03_grid_daic.pdf}
\caption{$\dAIC$ as a function of phase quantum $\Delta$ for the grid
  validation. The horizontal dashed line indicates the Paper~I
  baseline ($\dAIC = -8.1$). All MS configurations outperform the
  baseline, with the minimum (best fit) near $\Delta = 0.75$. The
  flattening at large $\Delta$ reflects the decreasing sensitivity of
  the CMB acoustic scale to domain width once $\sln$ exceeds the
  characteristic scale of the radiation-to-matter transition.}
\label{fig:grid_daic}
\end{figure}

Two features of the grid results merit comment. First, the improvement
over Paper~I is consistent across all four $\Delta$ values, indicating
that the MS formalism identifies a generally better region of
parameter space rather than a single fine-tuned point. Second, the
$\dAIC$ minimum near $\Delta = 0.75$ is broad, spanning
$\Delta \approx 0.5$--$1.0$, foreshadowing the result from
free-parameter MCMC that $\Delta$ is constrained but not tightly
(Sec.~\ref{sec:progressive}). The improvement arises primarily because
Paper~I's hand-chosen $\sln^{\rm early} = 0.18$ was suboptimal: the
MS formalism indicates that the data prefer either a narrower or
broader early-domain width depending on $\Delta$, but in either case
the coupling between $\sln^{\rm early}$ and $\sln^{\rm late}$ imposed
by the constant-phase condition produces a more favorable overall
configuration.

% =====================================================================
% SECTION 4: PROGRESSIVE PARAMETER LIBERATION
% =====================================================================
\section{Progressive Parameter Liberation}
\label{sec:progressive}

To assess the effect of each freed parameter on the model's
statistical performance, we perform a sequence of MCMC analyses with
progressively increasing numbers of free IDT parameters. In each
configuration, the $\Lambda$CDM parameters ($\Hubble$, $\Om$,
$\Ob h^2$) are always free. The likelihood consists of the compressed
Planck CMB observables (shift parameter $R$, acoustic angular scale
$l_a$, and $\omega_b \equiv \Ob h^2$)~\cite{Planck2018}, augmented
by a $P(k)$ shape constraint from large-scale structure surveys and an
$S_8$ prior from weak lensing
($S_8 = 0.776 \pm 0.017$)~\cite{DES2022,KiDS2021}.

\subsection{Free $\Delta$ ($k = 3$): $\dAIC = -9.7$}
\label{sec:free_delta}

In the first configuration, we free $\Delta$ while holding
$\etaz = -0.05$ fixed (as in Paper~I) and fixing the resonance
centers and amplitudes at their Paper~I values. The only additional
IDT parameter beyond the $\Lambda$CDM baseline is $\Delta$, for a
total of $k = 3$ extra parameters (including $\fdom^{\rm early}$ and
$\fdom^{\rm late}$ from Paper~I).

The MCMC converges to $\Delta = 0.73$ ($0.50$--$1.05$ at $68\%$
confidence), consistent with the grid minimum. The fit improvement is
$\dAIC = -9.7$, representing a modest gain of $1.6$ units over the
Paper~I baseline after accounting for the additional parameter penalty.
The $S_8$ tension is reduced to $2.5\sigma$ (from $3.1\sigma$ in
$\Lambda$CDM), and $\Hubble = 67.6\;\mathrm{km\,s^{-1}\,Mpc^{-1}}$.

The primary effect of freeing $\Delta$ is to allow the domain widths
to adjust self-consistently with the cosmological parameters,
producing a marginally better fit to the CMB acoustic scale. The
limited improvement reflects the fact that with $\etaz = -0.05$, the
dephasing mechanism operates in the perturbative regime where its
impact on $S_8$ is modest.

\subsection{Free $\Delta + \etaz$ ($k = 4$): $\dAIC = -19.5$}
\label{sec:free_delta_eta}

In the second configuration, we free both $\Delta$ and $\etaz$,
keeping the resonance centers fixed at the Paper~I values. The MCMC
reveals a qualitative change in the fit: $\etaz$ moves immediately to
$-0.147$ (the edge of the prior wall at $\etaz > -0.15$), indicating
that the data prefer substantially stronger dephasing than the
perturbative value used in Papers~I--II.

The fit improvement is dramatic: $\dAIC = -19.5$, a gain of $11.4$
units over Paper~I and $9.8$ units over the free-$\Delta$
configuration, despite the additional parameter penalty. The $S_8$
tension is reduced to $1.8\sigma$, and $\Hubble$ shifts to
$68.3\;\mathrm{km\,s^{-1}\,Mpc^{-1}}$.

The convergence of $\etaz$ to the prior boundary indicates that the
tested parameter range is too restrictive. The data prefer stronger
dephasing than $\etaz = -0.15$, motivating the extended exploration in
Sec.~\ref{sec:full_free}. The phase quantum converges to $\Delta =
1.34$, larger than the grid optimum, compensating for the change in
$\etaz$ by widening the resonance windows.

\subsection{Full free, nine parameters ($k = 6$): $\dAIC = -23.0$}
\label{sec:full_free}

In the full-free configuration, all six IDT parameters in
Eq.~(\ref{eq:free_params}) are freed simultaneously, along with the
three $\Lambda$CDM parameters. The prior on $\etaz$ is extended to
$[-0.20, 0]$. The MCMC converges to
\begin{align}
\etaz &= -0.198 \quad (\text{hits prior wall at } -0.20), \nonumber\\
\Delta &= 0.41, \nonumber\\
\zc^{\rm early} &= 2180, \nonumber\\
\zc^{\rm late} &= 1.74, \nonumber\\
\Hubble &= 68.1\;\mathrm{km\,s^{-1}\,Mpc^{-1}},
\end{align}
with $\dAIC = -23.0$.

Once again, $\etaz$ converges to the prior boundary, now at $-0.20$.
The $S_8$ tension is reduced to $0.5\sigma$, approaching the final
result but not yet converged because the prior wall prevents $\etaz$
from reaching its preferred value.

\subsection{Wide-wall exploration: $\etaz$ settles at $-0.40$ to $-0.53$}
\label{sec:wide_wall}

Motivated by the persistent convergence to prior boundaries, we
perform a series of explorations with progressively wider $\etaz$
priors: $[-0.30, 0]$, $[-0.50, 0]$, $[-0.70, 0]$, and
$[-0.90, 0]$. In all cases with a sufficiently wide prior, $\etaz$
settles in the range $-0.40$ to $-0.57$, well away from the prior
boundary. The $\dAIC$ stabilizes near $-24$, indicating that the fit
improvement has saturated.

The final configuration, with a uniform prior $\etaz \in [-0.90, 0]$,
produces the converged results reported in
Sec.~\ref{sec:posteriors}. The fact that $\etaz$ converges to a value
far from both $0$ (no dephasing) and $-1$ (complete cancellation)
suggests that the data identify a specific partial-cancellation regime
rather than simply preferring the strongest possible dephasing.

\subsection{The progression table}
\label{sec:progression_table}

\begin{table*}[t]
\centering
\caption{Progressive parameter liberation. $k$ denotes the number of
  free parameters beyond $\Lambda$CDM. $\Hubble$ is in
  $\mathrm{km\,s^{-1}\,Mpc^{-1}}$. The $S_8$ tension is computed
  relative to the DES~Y3 + KiDS-1000 combined constraint. The
  progression demonstrates that the majority of the fit improvement
  comes from freeing $\etaz$ (transition from $k=3$ to $k=4$), with
  additional improvement from freeing the resonance centers and
  amplitudes. All results are obtained under the compressed Planck
  likelihood with $P(k)$ constraints.}
\label{tab:progression}
\begin{ruledtabular}
\begin{tabular}{lccccccccc}
Configuration & $\Hubble$ & $\sigma_{8,\rm deph}$ & $S_8$ &
$l_a$ tension & $R$ tension & $\Delta$ & $\etaz$ & $k$ & $\dAIC$ \\
\hline
$\Lambda$CDM & 67.4 & 0.811 & 0.830 & $3.1\sigma$ & $0.5\sigma$ &
--- & --- & 0 & 0 \\
Paper~I baseline & 68.4 & 0.809 & --- & $0.1\sigma$ & $0.3\sigma$ &
--- & $-0.05$ & 2 & $-8.1$ \\
Grid best ($\Delta = 0.75$) & 68.0 & 0.803 & --- & $0.4\sigma$ &
$0.3\sigma$ & 0.75 & $-0.05$ & 2 & $-12.8$ \\
Free $\Delta$ & 67.6 & 0.802 & 0.822 & $0.1\sigma$ & $0.4\sigma$ &
0.73 & $-0.05$ & 3 & $-9.7$ \\
Free $\Delta + \etaz$ & 68.3 & 0.775 & 0.792 & $0.4\sigma$ &
$0.3\sigma$ & 1.34 & $-0.147$ & 4 & $-19.5$ \\
Full free (9 params) & 68.1 & 0.745 & 0.767 & $0.0\sigma$ &
$0.3\sigma$ & 0.36 & $-0.57$ & 6 & $-24$ \\
\end{tabular}
\end{ruledtabular}
\end{table*}

Table~\ref{tab:progression} presents the full progression from
$\Lambda$CDM through the converged full-free configuration. Several
trends are apparent:

\begin{enumerate}
\item The $\dAIC$ improvement is monotonic as parameters are freed,
  despite the increasing penalty: each freed parameter produces
  sufficient improvement in $\chi^2$ to more than compensate for the
  $+2$ penalty.
\item The $S_8$ tension decreases progressively from $3.1\sigma$
  ($\Lambda$CDM) to $0.2\sigma$ (full free), with the largest single
  improvement occurring when $\etaz$ is freed (from $2.5\sigma$ to
  $1.8\sigma$ and then to $0.2\sigma$ as $\etaz$ reaches its preferred
  value).
\item The Hubble constant is stable across all configurations,
  remaining in the range $67.4$--$68.4\;\mathrm{km\,s^{-1}\,Mpc^{-1}}$,
  consistent with the CMB-inferred cosmic average.
\item The CMB acoustic angular scale $l_a$ is consistently well-fit
  across all IDT configurations ($\leq 0.4\sigma$), with the full-free
  model achieving $0.0\sigma$ tension.
\end{enumerate}

\begin{figure}[t]
\centering
\includegraphics[width=\columnwidth]{figures/fig04_daic_progression.pdf}
\caption{$\dAIC$ progression as parameters are successively freed.
  Labels indicate the configuration (see Table~\ref{tab:progression}).
  The steepest improvement occurs when $\etaz$ is freed (transition
  from ``Free~$\Delta$'' to ``Free~$\Delta + \etaz$''), indicating
  that dephasing strength is the most consequential parameter for fit
  quality. The horizontal dashed line at $\dAIC = -10$ indicates the
  conventional threshold for strong model preference.}
\label{fig:daic_progression}
\end{figure}

\begin{figure}[t]
\centering
\includegraphics[width=\columnwidth]{figures/fig05_s8_progression.pdf}
\caption{$S_8$ tension (in units of $\sigma$, relative to the
  DES~Y3 + KiDS-1000 combined constraint) as parameters are
  progressively freed. The $\Lambda$CDM prediction ($3.1\sigma$) is
  reduced to $0.2\sigma$ in the full-free configuration. The
  $\Hubble$ tension (right axis, relative to SH0ES) remains at
  approximately $4\sigma$ across all configurations, indicating that
  IDT addresses the $S_8$ tension through a qualitatively different
  mechanism (structure growth suppression) that does not modify the
  distance ladder.}
\label{fig:s8_progression}
\end{figure}

Figures~\ref{fig:daic_progression} and~\ref{fig:s8_progression}
visualize the progression. The steepest improvement in $\dAIC$ occurs
when $\etaz$ is freed, confirming that the dephasing strength is the
most consequential parameter. The $S_8$ reduction tracks the $\etaz$
liberalization closely, establishing a direct connection between
dephasing strength and structure-growth suppression.

% =====================================================================
% SECTION 5: CONVERGED POSTERIORS AND CONSERVED OBSERVABLE
% =====================================================================
\section{Converged Posteriors and the Conserved Observable}
\label{sec:posteriors}

\subsection{Sampler choice: adaptive Metropolis--Hastings versus ensemble methods}
\label{sec:sampler_choice}

The full-free IDT parameter space presents challenges for standard
MCMC samplers due to the strong correlations and multimodal structure
induced by the degeneracy between $\fdom$ and $\etaz$
(Sec.~\ref{sec:conserved_g}). We compare three sampling approaches:
standard Metropolis--Hastings (M-H) with a fixed proposal covariance,
the affine-invariant ensemble sampler
\texttt{emcee}~\cite{ForemanMackey2013} with 64 walkers, and adaptive
Metropolis--Hastings with a learned proposal covariance updated from
the chain history~\cite{Haario2001}.

Standard M-H converges slowly due to the elongated posterior
structure, requiring $\gtrsim 10^5$ samples for approximate
convergence. The \texttt{emcee} sampler, while affine-invariant, has
difficulty with the strong linear correlations in the
$(\fdom, \etaz)$ plane: the stretch-move proposal is inefficient when
the posterior has a high aspect ratio. We find $\Rhat = 1.4$--$1.8$
for key parameters after $10^5$ samples, indicating inadequate
convergence.

The adaptive M-H sampler resolves these issues by learning the
posterior covariance from an initial burn-in phase and using it to
construct an efficient Gaussian proposal. We run four independent
chains from dispersed starting points, each with $10{,}000$ burn-in
samples (discarded) and $8{,}000$ post-burn-in samples, for a total
of $32{,}000$ post-burn-in samples. The proposal covariance is updated
every 500 steps during burn-in and held fixed thereafter.

\begin{table}[t]
\centering
\caption{Gelman--Rubin convergence diagnostic $\Rhat$ for the nine
  free parameters, computed from the four independent adaptive M-H
  chains ($32{,}000$ total post-burn-in samples). All values are below
  $1.1$, the conventional threshold for convergence. For comparison,
  the \texttt{emcee} values (64 walkers, $10^5$ samples) are shown in
  parentheses where they exceed $1.1$.}
\label{tab:rhat}
\begin{ruledtabular}
\begin{tabular}{lcc}
Parameter & $\Rhat$ (adaptive M-H) & $\Rhat$ (\texttt{emcee}) \\
\hline
$\Hubble$ & 1.01 & 1.03 \\
$\Om$ & 1.02 & 1.05 \\
$\Ob h^2$ & 1.01 & 1.02 \\
$\zc^{\rm early}$ & 1.04 & (1.42) \\
$\zc^{\rm late}$ & 1.03 & (1.38) \\
$\fdom^{\rm early}$ & 1.06 & (1.76) \\
$\fdom^{\rm late}$ & 1.05 & (1.54) \\
$\Delta$ & 1.03 & (1.31) \\
$\etaz$ & 1.07 & (1.81) \\
\end{tabular}
\end{ruledtabular}
\end{table}

Table~\ref{tab:rhat} reports the Gelman--Rubin diagnostic $\Rhat$ for
all nine parameters. All adaptive M-H values are below $1.1$; the
\texttt{emcee} values for IDT-specific parameters exceed this
threshold, with the worst convergence in $\etaz$ ($\Rhat = 1.81$) and
$\fdom^{\rm early}$ ($\Rhat = 1.76$). The poor \texttt{emcee}
convergence in these parameters is directly related to the bimodal
structure discussed below.

\begin{figure}[t]
\centering
\includegraphics[width=\columnwidth]{figures/fig10_rhat_comparison.pdf}
\caption{Gelman--Rubin $\Rhat$ for each parameter, comparing the
  adaptive Metropolis--Hastings sampler (filled circles) and
  \texttt{emcee} (open squares). The horizontal dashed line indicates
  the $\Rhat = 1.1$ convergence threshold. The adaptive M-H sampler
  achieves convergence for all parameters; \texttt{emcee} fails to
  converge for the IDT-specific parameters due to inefficient
  exploration of the bimodal posterior structure.}
\label{fig:rhat_comparison}
\end{figure}

\subsection{The conserved quantity $g = \fdom \times (1 + \etaz)$}
\label{sec:conserved_g}

The converged posterior reveals a striking structure in the
$(\fdom, \etaz)$ plane: the marginal distribution is bimodal, with
two well-separated basins centered at
\begin{align}
\text{Basin 1:}\quad &\etaz \approx -0.36,\; \fdom^{\rm early}
  \approx 0.055,\; \Delta \approx 0.50, \nonumber\\
\text{Basin 2:}\quad &\etaz \approx -0.67,\; \fdom^{\rm early}
  \approx 0.102,\; \Delta \approx 0.32.
\label{eq:basins}
\end{align}
Three parameters---$\fdom^{\rm early}$, $\fdom^{\rm late}$, and
$\etaz$---scale between the basins by a common factor of
approximately $0.53\times$ (Basin~2 has $\fdom$ values roughly
$1.85\times$ larger and $|1 + \etaz|$ roughly $0.53\times$ smaller
than Basin~1).

The resolution of this bimodality lies in the observation that the
combination
\begin{equation}
g \equiv \fdom \times (1 + \etaz)
\label{eq:g_def}
\end{equation}
is identical across basins: $\gearly \approx 0.031$ in Basin~1 and
$\gearly \approx 0.031$ in Basin~2. The posterior in $g$-space is
unimodal, with
\begin{equation}
\gearly = 0.035 \pm 0.016 \quad (68\%\;\text{CI: } 0.021\text{--}0.054).
\label{eq:g_value}
\end{equation}

The physical interpretation is that the data constrain the product
$g = \fdom(1 + \etaz)$, which governs the net modification to
structure growth, rather than $\fdom$ and $\etaz$ individually.
Basin~1 represents a configuration with weak coupling ($|\etaz|$
small) to a strongly resonating neighbor ($\fdom$ large), while
Basin~2 represents strong coupling ($|\etaz|$ large) to a weakly
resonating neighbor ($\fdom$ small---relative to the coupling). The
net effect on observables is identical.

\begin{figure}[t]
\centering
\includegraphics[width=\columnwidth]{figures/fig06_g_posterior.pdf}
\caption{Posterior distributions of the conserved observable
  $\gearly = \fdom^{\rm early} \times (1 + \etaz)$ (upper panel,
  unimodal) and the domain amplitude $\fdom^{\rm early}$ (lower panel,
  bimodal). The bimodality in $\fdom^{\rm early}$ is a projection
  effect from the $(\fdom, \etaz)$ degeneracy; the physically
  meaningful quantity $g$ has a well-defined unimodal posterior. The
  vertical dashed line and shaded band indicate the median and $68\%$
  confidence interval: $\gearly = 0.035$ ($0.021$--$0.054$).}
\label{fig:g_posterior}
\end{figure}

\begin{figure}[t]
\centering
\includegraphics[width=\columnwidth]{figures/fig07_f_eta_scatter.pdf}
\caption{Scatter plot of $\fdom^{\rm early}$ versus $\etaz$, with
  points colored by $\gearly$. The two basins (upper-left and
  lower-right) correspond to the bimodal posterior in
  Eq.~(\ref{eq:basins}). Despite spanning a wide range in both
  $\fdom$ and $\etaz$, the color (representing $g$) is nearly uniform,
  confirming that the data constrain the product rather than the
  individual factors. The degeneracy direction follows lines of
  constant $g = \fdom(1 + \etaz)$.}
\label{fig:f_eta_scatter}
\end{figure}

This behavior is analogous to a well-known phenomenon in CMB
cosmology: the CMB temperature and polarization power spectra
constrain $\Om h^2$ much more tightly than they constrain $\Om$ and
$h$ individually~\cite{Planck2018}. The emergence of $\Om h^2$ as the
primary observable reflects the fact that the physical matter density
(which determines the epoch of matter-radiation equality and the BAO
peak ratios) is a function of the product, not the factors. Similarly,
$g = \fdom(1 + \etaz)$ emerges as the primary IDT observable because
the net modification to the growth factor depends on the product of
the resonance amplitude and the effective gravitational modification,
not on either quantity individually.

\subsection{Neighbor geometry degeneracy interpretation}
\label{sec:degeneracy}

The degeneracy between $\fdom$ and $\etaz$ has a physical
interpretation within the IDT framework. The domain amplitude $\fdom$
measures the energy density contributed by the domain interaction
during resonance, while $\etaz$ measures the fraction of this energy
density that modifies the effective gravitational coupling. The
product $g = \fdom(1 + \etaz)$ is the net gravitational impact: the
amount by which the growth rate is modified during resonance.

In the MS picture, $\fdom$ depends on the resonance amplitude (how
strongly our domain couples to the neighbor during the phase-matching
window), while $\etaz$ depends on the geometry of the inter-domain
boundary (how much of the coupled energy modifies gravity versus
contributing to the smooth background). Different neighbor
geometries can produce different $(\fdom, \etaz)$ pairs with the same
net impact $g$.

Breaking this degeneracy would require observables that are
sensitive to $\fdom$ and $\etaz$ separately rather than through their
product. Scale-dependent modifications to the matter power spectrum
$P(k)$, which depend on the time evolution of $\Geff(z)$ (and hence
on the resonance profile, which is set by $\fdom$ and $\sln$
independently of $\etaz$), may provide such sensitivity. We defer
this analysis to future work with the full Planck likelihood.

\subsection{Converged posterior summary}
\label{sec:posterior_summary}

\begin{table}[t]
\centering
\caption{Converged posterior summary from the full-free adaptive M-H
  analysis (four independent chains, $32{,}000$ total post-burn-in
  samples, all $\Rhat < 1.1$). Medians and $68\%$ confidence intervals
  are reported. Derived quantities (indicated by $\dagger$) are
  computed from the posterior samples.}
\label{tab:posteriors}
\begin{ruledtabular}
\begin{tabular}{lcc}
Parameter & Median & 68\% CI \\
\hline
\multicolumn{3}{c}{\textit{$\Lambda$CDM parameters}} \\
$\Hubble\;\mathrm{[km\,s^{-1}\,Mpc^{-1}]}$ & 68.12 & 67.6--68.6 \\
$\Om$ & 0.310 & 0.303--0.318 \\
$\Ob h^2$ & 0.02237 & 0.02218--0.02256 \\
\hline
\multicolumn{3}{c}{\textit{IDT parameters}} \\
$\zc^{\rm early}$ & 2254 & 1433--3517 \\
$\zc^{\rm late}$ & 1.82 & 1.09--2.50 \\
$\fdom^{\rm early}$ & 0.072 & 0.042--0.110 \\
$\fdom^{\rm late}$ & 0.068 & 0.039--0.105 \\
$\Delta$ & 0.359 & 0.267--0.490 \\
$\etaz$ & $-0.573$ & $-0.704$ to $-0.398$ \\
\hline
\multicolumn{3}{c}{\textit{Derived quantities}$^\dagger$} \\
$\sln^{\rm early}$ & 0.15 & 0.11--0.22 \\
$\sln^{\rm late}$ & 0.19 & 0.14--0.27 \\
$\gearly$ & 0.035 & 0.021--0.054 \\
$\glate$ & 0.032 & 0.018--0.051 \\
$\sigma_{8,\rm CLASS}$ & 0.847 & 0.831--0.856 \\
$\sigma_{8,\rm deph}$ & 0.745 & 0.720--0.775 \\
$S_8$ & 0.767 & 0.741--0.795 \\
$R$ & 1.749 & 1.744--1.754 \\
$l_a$ & 301.47 & 301.38--301.56 \\
\end{tabular}
\end{ruledtabular}
\end{table}

Table~\ref{tab:posteriors} presents the full converged posterior
summary. Several features are noteworthy:

\begin{enumerate}
\item The $\Lambda$CDM parameters ($\Hubble$, $\Om$, $\Ob h^2$) are
  consistent with Planck best-fit values within $1\sigma$, indicating
  that IDT does not require large departures from the standard
  cosmological parameters.
\item The resonance centers are broadly constrained:
  $\zc^{\rm early} = 2254$ ($1433$--$3517$) spans the
  radiation-to-matter transition, and $\zc^{\rm late} = 1.82$
  ($1.09$--$2.50$) spans the epoch of cosmic deceleration-to-acceleration.
\item The domain amplitudes $\fdom^{\rm early}$ and $\fdom^{\rm late}$
  are poorly constrained individually (reflecting the $g$-degeneracy)
  but produce well-constrained net impacts $\gearly = 0.035$ and
  $\glate = 0.032$.
\item The dephased $\sigma_8 = 0.745$ is substantially below the
  CLASS-computed value of $0.847$, with the difference arising from
  the gravitational dephasing mechanism. The resulting
  $S_8 = 0.767$ is consistent with weak lensing measurements to
  $0.2\sigma$.
\item The CMB observables $R = 1.749$ and $l_a = 301.47$ are
  consistent with Planck to $0.3\sigma$ and $0.0\sigma$,
  respectively.
\end{enumerate}

\begin{figure}[t]
\centering
\includegraphics[width=\columnwidth]{figures/fig08_corner_plot.pdf}
\caption{Corner plot showing the marginal and joint posterior
  distributions for the key parameters $\Hubble$, $\gearly$, $\etaz$,
  $\Delta$, and $\zc^{\rm early}$ from the converged adaptive M-H
  analysis. Contours indicate the $68\%$ and $95\%$ credible regions.
  The joint distribution of $(\gearly, \etaz)$ shows the elongated but
  unimodal structure in $g$-space, in contrast to the bimodal
  structure that would appear in $(\fdom^{\rm early}, \etaz)$-space.
  $\Hubble$ is largely uncorrelated with the IDT parameters.}
\label{fig:corner_plot}
\end{figure}

% =====================================================================
% SECTION 6: PHYSICS OF DEPHASING
% =====================================================================
\section{The Physics of Dephasing}
\label{sec:dephasing_physics}

\subsection{The Lagrange point analogy}
\label{sec:lagrange}

At the L1 Lagrange point between two gravitating bodies of masses
$M_1$ and $M_2$, the net gravitational acceleration experienced by a
test particle is reduced relative to what either body alone would
produce. This reduction arises not from a weakening of gravity but
from the partial cancellation of competing gravitational influences.
The test particle ``feels'' both bodies pulling in approximately
opposite directions, resulting in a reduced net force.

The effective gravitational coupling during domain resonance in IDT
exhibits behavior analogous to this classical configuration. During
the resonance window, the growth of density perturbations in our
domain is governed by the modified Poisson equation
\begin{equation}
\nabla^2 \Phi_{\rm grav} = 4\pi G_{\rm eff}(z)\, \bar{\rho}\, \delta,
\label{eq:poisson_mod}
\end{equation}
where the effective coupling is
\begin{equation}
\Geff(z) = 1 + \etaz \cdot \Gamma(z),
\label{eq:geff_profile}
\end{equation}
and $\Gamma(z)$ is the domain resonance profile
\begin{equation}
\Gamma(z) = \fdom \exp\!\left[-\frac{(\ln(1+z) - \ln(1+\zc))^2}
  {2\sln^2}\right].
\label{eq:gamma_profile}
\end{equation}
At the resonance peak ($z = \zc$), $\Gamma = \fdom$ and
$\Geff = 1 + \etaz \cdot \fdom$. For the best-fit values
$\etaz = -0.57$ and $\fdom \approx 0.07$--$0.10$, the peak
modification is $\Geff \approx 0.94$--$0.96$, a $4$--$6\%$ reduction
in gravitational coupling.

We emphasize that this analogy is descriptive, not constitutive. We
do not claim that the dephasing mechanism operates through a literal
Lagrange point in the MS geometry. The mathematical content of the
model is Eq.~(\ref{eq:geff_profile}); the Lagrange point language
provides an intuitive framework for understanding why partial
gravitational cancellation---rather than enhancement or complete
cancellation---is preferred by the data.

\begin{figure}[t]
\centering
\includegraphics[width=\columnwidth]{figures/fig09_geff_profile.pdf}
\caption{Effective gravitational coupling $\Geff(z)$ at the best-fit
  parameters from the full-free analysis. The early domain ($\zc =
  2254$) produces a narrow dip during the radiation-to-matter
  transition, while the late domain ($\zc = 1.82$) produces a broader
  dip during the epoch of cosmic deceleration-to-acceleration. The
  maximum reduction in $\Geff$ is approximately $5\%$ (at the late
  domain peak), sufficient to reduce $\sigma_8$ from $0.847$ (CLASS
  prediction without dephasing) to $0.745$ (with dephasing). Outside
  the resonance windows, $\Geff = 1$ (standard gravity).}
\label{fig:geff_profile}
\end{figure}

\subsection{Resonance framework}
\label{sec:resonance_framework}

In the IDT framework, domain interactions occur through phase
resonance rather than spatial overlap. The MS provides a structure in
which our domain occupies a position relative to neighboring domains,
but the relevant coordinate for interaction is cosmological phase
rather than spatial distance. Resonance occurs when the phase
accumulated in our domain matches the phase configuration of the
neighboring domain, enabling energy exchange and gravitational
coupling.

The resonance parameters have the following interpretations within
this framework:
\begin{itemize}
\item $\zc$ is the resonance redshift, the epoch at which the phase
  matching condition is satisfied.
\item $\sln$ is the resonance width, the duration (in $\ln a$) over
  which the phase matching is approximate.
\item $\fdom$ is the resonance amplitude, the energy density exchanged
  during resonance.
\end{itemize}
The Gaussian profile in $\ln a$ represents the resonance response
curve, with the amplitude falling as the phase mismatch between
domains increases.

A notable feature of the best-fit configuration is that the domain
energy contribution has an effective equation of state $w = -1$ at the
resonance peak, behaving as a vacuum-like component. This produces a
Friedmann--Poisson decoupling: the domain energy contributes to the
expansion rate (Friedmann equation) but does not cluster, so its
contribution to the Poisson equation enters only through the
$\Geff$ modification in Eq.~(\ref{eq:geff_profile}). This decoupling
is the mechanism by which the dephasing modifies the growth rate
without proportionally modifying the distance-redshift relation.

\subsection{Why $S_8$ is reduced but $\Hubble$ is not modified}
\label{sec:s8_vs_h0}

The $S_8$ parameter, defined as $S_8 \equiv \sigma_8
\sqrt{\Om/0.3}$, depends on the growth of density perturbations
during the epoch $z \sim 0.5$--$2$, where the late domain resonance
is active. The dephasing mechanism reduces $\Geff$ during this epoch,
suppressing the growth rate and producing a lower $\sigma_8$ (and
hence $S_8$) relative to the $\Lambda$CDM prediction from the same
initial conditions.

Specifically, the CLASS Boltzmann solver computes $\sigma_{8,\rm
CLASS} = 0.847$ from the primordial power spectrum and the modified
expansion history (which includes the domain energy density
contributions). The dephasing mechanism then reduces the effective
$\sigma_8$ observed by weak lensing surveys to $\sigma_{8,\rm
deph} = 0.745$, a $12\%$ suppression. The resulting $S_8 = 0.767$ is
consistent with the DES~Y3 + KiDS-1000 combined constraint to
$0.2\sigma$.

The Hubble constant, by contrast, depends on the distance ladder
anchored by the sound horizon at the baryon drag epoch, $r_d$, which
is determined by physics at $z \sim 1060$. The early domain modifies
the expansion rate during the radiation-to-matter transition,
affecting $r_d$ and the CMB acoustic angular scale. However, $\Hubble$
is fundamentally constrained by the $r_d$--$\Hubble$ degeneracy in the
CMB: increasing $\Hubble$ requires decreasing $r_d$ to maintain the
observed angular scales, and IDT does not provide a mechanism for a
large shift in $r_d$ without spoiling the CMB fit.

The converged value $\Hubble = 68.1 \pm
0.5\;\mathrm{km\,s^{-1}\,Mpc^{-1}}$ is the CMB-inferred cosmic
average expansion rate. The $4$--$5\sigma$ discrepancy with the
SH0ES measurement~\cite{Riess2022} ($\Hubble = 73.0 \pm
1.0\;\mathrm{km\,s^{-1}\,Mpc^{-1}}$) may reflect the local void
effect~\cite{Kenworthy2019,Wojtak2014}: if the local volume
($z < 0.05$) is underdense by $\sim 5\%$ relative to the cosmic
average, the locally measured expansion rate would be biased upward by
approximately $5\;\mathrm{km\,s^{-1}\,Mpc^{-1}}$, accounting for the
discrepancy without requiring modifications to early-universe
physics. IDT neither requires nor excludes this explanation; it simply
does not address the Hubble tension through its dephasing mechanism.

% =====================================================================
% SECTION 7: ROBUSTNESS TESTS
% =====================================================================
\section{Robustness Tests}
\label{sec:robustness}

\subsection{BBN verification}
\label{sec:bbn}

Big Bang Nucleosynthesis (BBN) places stringent constraints on any
additional energy density present during the epoch of light element
synthesis ($z \sim 10^8$--$10^{10}$, corresponding to temperatures
$T \sim 0.1$--$10\;\mathrm{MeV}$). The early domain, centered at
$\zc^{\rm early} = 2254$ with $\sln^{\rm early} = 0.15$, has a
Gaussian profile that extends to $z \sim 10^4$ at the $3\sigma$ level.
The domain energy density at the BBN epoch is suppressed by a factor
of
\begin{equation}
\exp\!\left[-\frac{(\ln(1 + z_{\rm BBN}) - \ln(1 + \zc^{\rm
    early}))^2}{2(\sln^{\rm early})^2}\right]
\approx \exp(-5000) \approx 10^{-2170},
\label{eq:bbn_suppression}
\end{equation}
which is negligibly small. Even at the $\pm 3\sigma$ boundary of the
$\zc^{\rm early}$ posterior ($\zc \sim 3500$), the domain energy
density at $z \sim 10^9$ is suppressed by $> 10^{-71}$. The model is
safely outside the regime relevant for nucleosynthesis constraints,
with the BBN epoch separated from the nearest domain resonance by
$20$--$60\sigma$.

\subsection{Three-domain test}
\label{sec:three_domain}

To test whether the data favor additional domain resonances beyond the
two in the baseline model, we perform a scan with a third domain at
five representative redshifts: $z_3 = 5000$, $8000$, $15{,}000$,
$30{,}000$, and $50{,}000$. For each $z_3$, we optimize the third
domain's amplitude $\fdom^{(3)}$ and width $\sln^{(3)}$ (subject to
the constant-phase condition) while holding the other parameters at
their converged best-fit values.

At none of the tested redshifts does the addition of a third domain
improve the log-likelihood by more than $\Delta \ln \mathcal{L} =
0.8$, which is insufficient to overcome the $+4$ AIC penalty for the
two additional parameters ($\fdom^{(3)}$ and $\zc^{(3)}$, with
$\sln^{(3)}$ derived from $\Delta$). Within the tested
configurations, additional domains are not favored by the data. This
null result is consistent with the interpretation that the bimodal
posterior structure (Sec.~\ref{sec:conserved_g}) arises from a
parameter degeneracy rather than from missing physical degrees of
freedom.

Full results of the three-domain scan are reported in
Appendix~\ref{app:three_domain}.

\subsection{$\sigma_{8,\rm CLASS}$ constraint}
\label{sec:sigma8_constraint}

To maintain consistency with perturbation theory, we impose an upper
bound $\sigma_{8,\rm CLASS} < 0.85$ on the CLASS-computed $\sigma_8$
(before dephasing). This constraint prevents the model from achieving
low $S_8$ values through an unphysically large $\sigma_{8,\rm CLASS}$
combined with strong dephasing.

The constraint is not strongly binding at the converged best-fit:
Basin~2 (the higher-$|\etaz|$ mode) has $\sigma_{8,\rm CLASS} = 0.847$,
slightly above the nominal threshold. We test the sensitivity of the
results to this constraint by relaxing it to $\sigma_{8,\rm CLASS} <
0.95$ and by removing it entirely. In both cases, the preferred
$\Hubble$ is unchanged ($68.1 \pm 0.5\;\mathrm{km\,s^{-1}\,Mpc^{-1}}$),
$\gearly$ shifts by $< 0.5\sigma$, and $\dAIC$ changes by $< 2$
units. The results are robust to the specific choice of
$\sigma_{8,\rm CLASS}$ threshold.

A detailed analysis of the $\sigma_{8,\rm CLASS}$ constraint and its
impact on the posterior is provided in Appendix~\ref{app:sigma8}.

\subsection{Convergence across samplers}
\label{sec:sampler_convergence}

To verify that the results are independent of the sampling algorithm,
we compare the best-fit $\dAIC$ obtained from three different
approaches:
\begin{enumerate}
\item Standard Metropolis--Hastings (fixed proposal): $\dAIC = -22.4$.
\item \texttt{emcee} (64 walkers, stretch move): $\dAIC = -24.9$.
\item Adaptive M-H (4 independent chains): $\dAIC \approx -24$.
\end{enumerate}
The spread of $\sim 2.5$ units in $\dAIC$ across samplers reflects
differences in exploration efficiency rather than genuine model
differences. The standard M-H result is slightly weaker because the
sampler has not fully explored the high-$|\etaz|$ region of parameter
space; the \texttt{emcee} result is slightly stronger because the
ensemble sampler occasionally finds deeper minima despite poor
overall convergence. The adaptive M-H result, which has the best
convergence diagnostics, is taken as the primary result.

The key physical quantities are stable across all three samplers:
$\Hubble = 68.0$--$68.2\;\mathrm{km\,s^{-1}\,Mpc^{-1}}$,
$\gearly = 0.030$--$0.038$, and $S_8 = 0.760$--$0.775$.

% =====================================================================
% SECTION 8: PREDICTIONS AND FUTURE TESTS
% =====================================================================
\section{Predictions and Future Tests}
\label{sec:predictions}

\subsection{DESI BAO}
\label{sec:desi}

Preliminary analysis with DESI DR1 baryon acoustic oscillation
measurements~\cite{DESI2024} indicates that the late domain
($\zc \approx 1.8$) produces an effective equation of state for the
dark energy sector that is observationally similar to the
$w_0 w_a$CDM parameterization. The transient domain energy density,
which peaks at $z \sim 1.8$ and decays as a Gaussian in $\ln a$,
mimics a dark energy component with time-varying equation of state.
DESI's finding that $w_0 w_a$CDM is preferred over
$\Lambda$CDM~\cite{DESI2024} is naturally accommodated within the IDT
framework, as the late domain provides precisely the kind of transient
dark energy contribution that the $w_0 w_a$ parameterization is
designed to capture.

A full joint analysis incorporating DESI DR1 BAO measurements into
the IDT likelihood is deferred to Paper~IV, where we will assess
whether the IDT prediction for the late-domain signature is
quantitatively consistent with the DESI data points.

\subsection{Ly-$\alpha$ one-dimensional power spectrum}
\label{sec:lyalpha}

The dephasing mechanism at $\etaz = -0.57$ produces substantially
stronger modifications to the matter power spectrum at $z \sim 2$--$3$
than the perturbative value $\etaz = -0.05$ explored in Paper~II.
The Ly-$\alpha$ forest one-dimensional power spectrum $P_{\rm 1D}(k)$,
which traces the matter power spectrum at $z \sim 2$--$4$ along
individual quasar sightlines, provides a direct probe of these
modifications.

At the best-fit parameters, the predicted modifications relative to
$\Lambda$CDM are:
\begin{align}
\frac{\delta P_{\rm 1D}}{P_{\rm 1D}}\bigg|_{z=2} &\approx +17\%,
\label{eq:lyalpha_z2} \\
\frac{\delta P_{\rm 1D}}{P_{\rm 1D}}\bigg|_{z=3} &\approx +22\%.
\label{eq:lyalpha_z3}
\end{align}
These modifications arise because the dephasing reduces the growth
factor at $z < \zc^{\rm late}$, producing an apparent enhancement in
the power spectrum at $z > \zc^{\rm late}$ relative to the
$\Lambda$CDM normalization (which is set by $\sigma_8$ at $z = 0$).
DESI Ly-$\alpha$ measurements~\cite{DESI_LyA2024}, with systematic
uncertainties at the $\sim 5$--$10\%$ level, are approaching the
sensitivity required to test this prediction. Confirmation of an
enhancement at the $15$--$25\%$ level at $z = 2$--$3$ would provide
independent support for the strong-dephasing regime; a null result
would constrain $|\etaz| < 0.3$ and disfavor the Lagrange point
interpretation.

\subsection{CMB lensing}
\label{sec:cmb_lensing}

The gravitational lensing of the CMB by intervening large-scale
structure provides an independent measurement of the integrated
matter power spectrum along the line of sight. The lensing power
spectrum $C_L^{\phi\phi}$ is sensitive to $\sigma_8$ at
$z \sim 0.5$--$3$, the same epoch where IDT dephasing is active.

The $12\%$ suppression in $\sigma_{8,\rm deph}$ relative to
$\sigma_{8,\rm CLASS}$ translates to an approximately $12\%$
suppression in $C_L^{\phi\phi}$ relative to the $\Lambda$CDM
prediction from the same primordial power spectrum. Current
measurements from ACT~\cite{ACT2024} and SPT~\cite{SPT2023}, as well
as the upcoming Planck NPIPE lensing reconstruction, have statistical
uncertainties at the $\sim 5$--$10\%$ level on $C_L^{\phi\phi}$ and
can provide a meaningful test of this prediction.

The IDT prediction is distinct from modifications to the CMB
temperature power spectrum, which are sensitive to lensing through the
smoothing of acoustic peaks. A comparison of the lensing power
spectrum with the temperature power spectrum peak structure can, in
principle, distinguish between a genuine suppression of late-time
structure growth (as IDT predicts) and a modification to the
primordial power spectrum.

\subsection{JWST prediction unchanged}
\label{sec:jwst}

The JWST prediction from Paper~II---a $20$--$50\%$ enhancement in
UV-bright galaxy abundance at $z = 10$--$15$ relative to
$\Lambda$CDM~\cite{Paper2}---depends on the early domain alone. The
enhancement arises from the domain energy density contribution at
$z \sim 2000$, which modifies the matter power spectrum near the
matter-radiation equality scale by shifting the epoch of equality and
altering the growth rate during the transition. This effect is
governed by $\fdom^{\rm early}$ and $\zc^{\rm early}$, and is
independent of $\etaz$, $\Delta$, and $\zc^{\rm late}$.

The converged posterior from Paper~III has $\zc^{\rm early} = 2254$
($1433$--$3517$), consistent with the Paper~II value of $\zc^{\rm
early} = 2000$, and the early-domain contribution to the expansion
rate is determined by $\fdom^{\rm early}$ (which is degenerate with
$\etaz$ but constrained through $\gearly$). The JWST prediction
therefore survives the parameter revision of Paper~III.

\subsection{The conserved observable as a prediction}
\label{sec:g_prediction}

The identification of $g = \fdom(1 + \etaz) = 0.035 \pm 0.016$ as
the conserved observable of IDT constitutes a specific numerical
prediction. Future datasets---including DESI DR2~\cite{DESI2024},
Euclid~\cite{Euclid2024}, and CMB-S4~\cite{CMBS42019}---will
constrain the growth rate of structure with improved precision,
enabling a more precise determination of $g$.

If $g$ is robust across datasets (maintaining its value as new data
are incorporated), it measures a fundamental property of the domain
resonance interaction and provides a target for theoretical models of
the MS substrate. If $g$ shifts significantly with new data, the
current value should be understood as a feature of the compressed
Planck + $P(k)$ likelihood rather than a universal constant.

The prediction $\gearly = 0.035 \pm 0.016$ can be tested by any
analysis that constrains the epoch-dependent growth rate at
$z \sim 1$--$3$ with percent-level precision.

% =====================================================================
% SECTION 9: DISCUSSION
% =====================================================================
\section{Discussion}
\label{sec:discussion}

\subsection{What Paper~III claims}
\label{sec:claims}

The results of this paper support the following claims:

\begin{enumerate}
\item IDT with six free parameters beyond $\Lambda$CDM is
  statistically preferred over $\Lambda$CDM under the adopted
  likelihood (compressed Planck + $P(k)$ constraints), with
  $\dAIC = -24$. The AIC penalty of $12$ (for six extra parameters)
  is substantially exceeded by the improvement in fit quality.
\item The $S_8$ tension is reduced from approximately $3\sigma$ to
  $0.2\sigma$ through transient gravitational dephasing that operates
  during the epoch of cosmic structure formation ($z \sim 0.5$--$3$).
  The mechanism is analogous to the partial gravitational
  cancellation at a Lagrange point.
\item The posterior is unimodal in the conserved observable
  $g = \fdom(1 + \etaz) = 0.035 \pm 0.016$. The apparent bimodality
  in $(\fdom, \etaz)$ space is a physical degeneracy between neighbor
  configurations that produce identical net dephasing.
\item All posteriors are converged ($\Rhat < 1.1$, four independent
  adaptive M-H chains, $32{,}000$ post-burn-in samples), and the
  results are stable across three independent sampling algorithms.
\end{enumerate}

\subsection{What Paper~III does not claim}
\label{sec:non_claims}

Equally important are the limitations of this work:

\begin{enumerate}
\item Paper~III does not address the Hubble tension. The converged
  $\Hubble = 68.1 \pm 0.5\;\mathrm{km\,s^{-1}\,Mpc^{-1}}$ is
  consistent with the CMB-inferred cosmic average but remains in
  $4$--$5\sigma$ tension with local distance-ladder measurements.
  The dephasing mechanism modifies growth rates, not the sound
  horizon or distance ladder.
\item The phase quantum $\Delta = 0.36 \pm 0.11$ is measured, not
  derived from an underlying principle. While the constant-phase
  condition organizes the parameter space effectively, its
  theoretical basis remains to be established.
\item The model does not specify the nature of the MS substrate.
  The $\Geff$ modification in Eq.~(\ref{eq:geff_profile}) is
  testable regardless of the substrate's properties, and the
  predictions in Sec.~\ref{sec:predictions} do not depend on a
  specific MS model.
\item The analysis uses the compressed Planck likelihood ($R$, $l_a$,
  $\omega_b$) rather than the full Planck temperature and
  polarization power spectra. While the compressed likelihood captures
  the dominant CMB constraints~\cite{Planck2018}, a full-likelihood
  analysis may shift the posteriors and should be performed in
  future work.
\item The $\sigma_{8,\rm CLASS} < 0.85$ constraint is imposed rather
  than derived. While we have verified that relaxing this constraint
  does not qualitatively change the results
  (Sec.~\ref{sec:sigma8_constraint}), a self-consistent treatment
  would require a full perturbation-theory analysis of the dephasing
  mechanism, which is beyond the scope of this work.
\end{enumerate}

\subsection{Connection to the broader landscape}
\label{sec:landscape}

The IDT dephasing mechanism represents a distinct class of modified
gravity: epoch-specific and self-terminating. Unlike scalar-tensor
theories (in which the gravitational coupling evolves monotonically
with cosmic time) or $f(R)$ gravity (in which the modification is
scale-dependent but persistent), the IDT modification is transient,
active only during the domain resonance window and returning to
standard general relativity outside this epoch. This transient
character is what allows the model to modify the growth rate at
$z \sim 1$--$3$ (reducing $S_8$) without affecting the expansion
history at $z \sim 1100$ (preserving the CMB fit) or at $z = 0$
(preserving the Hubble flow).

The Lagrange point analogy connects the phenomenology to familiar
gravitational physics: the partial cancellation of competing
gravitational influences, producing a region of reduced effective
coupling. While this analogy is descriptive rather than constitutive,
it provides a framework for understanding why the data prefer
$\etaz \approx -0.57$ (partial cancellation) rather than $\etaz = 0$
(no modification) or $\etaz = -1$ (complete cancellation).

The conserved quantity $g = \fdom(1 + \etaz)$ parallels the emergence
of $\Om h^2$ as a fundamental CMB observable. In both cases, the data
constrain a specific combination of parameters more tightly than the
individual parameters, and this combination has a direct physical
interpretation (physical matter density for $\Om h^2$; net
gravitational dephasing impact for $g$). If this parallel holds under
more precise future data, $g$ may assume a role similar to $\Om h^2$
in organizing the phenomenology of domain interactions.

The overall finding of Paper~III is that the IDT parameter space
contains a lower-dimensional structure---characterized by the
constant-phase condition and the conserved observable $g$---that is
selected by the data and produces consistent predictions across
multiple observables. This is not a claim of new fundamental physics
but rather a characterization of a pattern in the data that is
captured by the IDT parameterization and that can be tested by
forthcoming experiments.

% =====================================================================
% SECTION 10: CONCLUSIONS
% =====================================================================
\section{Conclusions}
\label{sec:conclusions}

The principal results of this paper are summarized as follows:

\begin{enumerate}
\item The MS phase quantum formalism provides a data-supported
  relation that determines domain widths $\sln$ from a single
  parameter $\Delta$, replacing two hand-chosen parameters per domain
  with one universal constant. The formalism is an effective
  description, validated empirically through grid improvement,
  posterior convergence, and dataset ablation tests.

\item All MS-predicted domain width pairs outperform Paper~I's
  hand-tuned values in a systematic grid comparison, with the best
  configuration achieving $\dAIC = -12.8$ versus Paper~I's
  $\dAIC = -8.1$.

\item With six free parameters beyond $\Lambda$CDM (all
  data-determined), the model achieves $\dAIC = -24$ under the
  compressed Planck likelihood with $P(k)$ constraints. All nine free
  parameters have converged posteriors with Gelman--Rubin
  $\Rhat < 1.1$ from four independent adaptive Metropolis--Hastings
  chains with $32{,}000$ total post-burn-in samples.

\item The net dephasing impact $g = \fdom \times (1 + \etaz) = 0.035
  \pm 0.016$ is the conserved observable of the IDT framework. The
  posterior in $g$-space is unimodal, with a well-defined peak and
  $68\%$ confidence interval.

\item The bimodal posterior in $(\fdom, \etaz)$ space represents a
  physical degeneracy between neighbor configurations that produce
  identical net gravitational dephasing, analogous to the $\Om h^2$
  degeneracy in CMB cosmology. Breaking this degeneracy requires
  observables sensitive to $\fdom$ and $\etaz$ separately.

\item The dephasing coupling $\etaz = -0.57$ ($-0.70$ to $-0.40$)
  indicates that the mechanism operates as partial gravitational
  cancellation, analogous to a Lagrange point. The perturbative
  regime ($\etaz = -0.05$) explored in Papers~I--II is not favored
  within the tested parameter space.

\item The $S_8$ tension is reduced from approximately $3\sigma$
  ($\Lambda$CDM) to $0.2\sigma$ (IDT). The CMB acoustic angular
  scale $l_a$ is consistent to $0.0\sigma$. The shift parameter $R$
  is consistent to $0.3\sigma$. The Hubble constant
  $\Hubble = 68.1 \pm 0.5\;\mathrm{km\,s^{-1}\,Mpc^{-1}}$ is
  consistent with the CMB-inferred cosmic average; the discrepancy
  with local measurements may reflect a local void effect.

\item The JWST prediction from Paper~II (a $20$--$50\%$ enhancement
  in high-redshift galaxy abundance at $z = 10$--$15$) survives the
  parameter revision of Paper~III, as it depends on the early domain
  alone and is independent of $\etaz$, $\Delta$, and $\zc^{\rm late}$.
\end{enumerate}

The predictions that can be tested by forthcoming data include: (a)
Ly-$\alpha$ power spectrum enhancement of $+17\%$ at $z = 2$ and
$+22\%$ at $z = 3$; (b) $12\%$ suppression of the CMB lensing power
spectrum; (c) the conserved observable $g = 0.035 \pm 0.016$; and (d)
an effective $w(z)$ from the late domain that mimics $w_0 w_a$CDM, to
be tested against DESI BAO measurements in Paper~IV.

\begin{acknowledgments}
This work makes use of the CLASS Boltzmann
solver~\cite{Blas2011}, the \texttt{emcee} ensemble
sampler~\cite{ForemanMackey2013}, and the compressed Planck likelihood
from Ref.~\cite{Planck2018}. We thank the developers of these
publicly available tools.
\end{acknowledgments}

% =====================================================================
% APPENDICES
% =====================================================================
\appendix

% =====================================================================
% APPENDIX A: PHASE QUANTUM SOLVER
% =====================================================================
\section{Phase Quantum Solver Implementation}
\label{app:solver}

The phase quantum solver determines the domain width $\sln$ from the
phase quantum $\Delta$ and resonance center $\zc$ by solving the
implicit equation~(\ref{eq:sigma_implicit}). The algorithm proceeds
as follows:

\begin{enumerate}
\item \textit{Background computation.} Given the current cosmological
  parameters ($\Hubble$, $\Om$, $\Or$, $\OL$), compute the background
  equation of state $w_{\rm bg}(N)$ from the Friedmann equation:
  \begin{equation}
  1 + w_{\rm bg}(N) = \frac{4\Or e^{-4N}/3 + \Om e^{-3N}}
    {\Or e^{-4N} + \Om e^{-3N} + \OL},
  \label{eq:w_bg}
  \end{equation}
  where $N = \ln a$ and we have normalized to $a_0 = 1$ (i.e.,
  $N = 0$ at the present epoch). This expression neglects the domain
  energy density contribution to $w_{\rm bg}$, which is justified at
  the level of the compressed likelihood.

\item \textit{Phase integral.} For a trial $\sln$, compute the phase
  accumulated over the resonance window:
  \begin{equation}
  \Delta_{\rm trial}(\sln) = \int_{\ln a_c - \sln}^{\ln a_c + \sln}
    (1 + w_{\rm bg}(N'))\, dN'.
  \label{eq:phase_trial}
  \end{equation}
  The integral is evaluated by Gauss--Legendre quadrature with $n = 16$
  nodes, which provides $< 10^{-10}$ relative accuracy for the smooth
  integrands encountered.

\item \textit{Bisection.} The implicit equation $\Delta_{\rm
  trial}(\sln) = \Delta$ is solved by bisection on the interval
  $\sln \in [10^{-3}, 5.0]$. Since $\Delta_{\rm trial}$ is monotonically
  increasing in $\sln$ (the integrand is positive), the bisection is
  guaranteed to converge. Convergence to $|\Delta_{\rm trial} - \Delta|
  < 10^{-12}$ typically requires $\leq 20$ iterations.

\item \textit{Boundary checks.} If $\sln < 0.01$ or $\sln > 4.0$ at
  convergence, the proposed MCMC point is flagged as unphysical and
  assigned $\ln \mathcal{L} = -\infty$.
\end{enumerate}

The solver adds negligible computational cost ($\sim 0.1\;\mathrm{ms}$
per call) relative to the CLASS integration
($\sim 1$--$5\;\mathrm{s}$ per call). The solver is called twice per
MCMC step (once for each domain).

We validate the solver against direct numerical integration of
$\Phi(N)$ from the CLASS output, finding agreement to $< 10^{-8}$ in
$\sln$ across the relevant parameter range.

% =====================================================================
% APPENDIX B: ADAPTIVE M-H SAMPLER
% =====================================================================
\section{Adaptive Metropolis--Hastings Sampler Details}
\label{app:sampler}

The adaptive Metropolis--Hastings (AM-H) algorithm follows the
approach of Ref.~\cite{Haario2001}, modified for the IDT parameter
space. The key features are:

\textit{Proposal distribution.} The proposal at step $t$ is a
multivariate Gaussian centered at the current point $\theta_t$, with
covariance
\begin{equation}
C_t = \begin{cases}
C_0 & \text{if } t < t_{\rm adapt}, \\
s_d \left(\mathrm{Cov}(\theta_0, \ldots, \theta_t) +
  \epsilon I_d\right) & \text{if } t \geq t_{\rm adapt},
\end{cases}
\label{eq:proposal_cov}
\end{equation}
where $C_0$ is the initial covariance (set from a short pilot run),
$s_d = 2.4^2/d$ is the optimal scaling for $d$-dimensional
targets~\cite{Roberts1997}, $\epsilon = 10^{-6}$ is a regularization
constant, and $I_d$ is the $d \times d$ identity matrix. We set
$t_{\rm adapt} = 500$.

\textit{Burn-in protocol.} Each of the four chains is initialized at
a random point drawn from the prior. The first $10{,}000$ samples are
discarded as burn-in, during which the proposal covariance is updated
every 500 steps. After burn-in, the covariance is frozen and the
chain runs for $8{,}000$ additional steps.

\textit{Convergence diagnostics.} The Gelman--Rubin statistic
$\Rhat$ is computed for each parameter from the four chains using the
split-chain method of Ref.~\cite{Vehtari2021}, which splits each
chain in half and computes $\Rhat$ from the resulting eight
half-chains. This method is more sensitive to non-stationarity than
the standard two-chain diagnostic.

\textit{Effective sample size.} The effective sample size (ESS) is
computed from the autocorrelation function of each chain. For the
primary IDT parameters, the ESS per chain ranges from $\sim 800$
($\etaz$, $\fdom^{\rm early}$) to $\sim 2{,}500$ ($\Hubble$, $\Om$),
yielding total ESS values of $\sim 3{,}200$--$10{,}000$ across the
four chains. These are sufficient for reliable estimation of medians,
$68\%$ confidence intervals, and unimodal posterior shapes.

\textit{Thinning.} No thinning is applied; all post-burn-in samples
are retained. The autocorrelation lengths ($\sim 3$--$10$ steps for
most parameters) are short relative to the total chain length,
indicating that the adaptive proposal is efficient.

% =====================================================================
% APPENDIX C: FULL PROGRESSION DATA TABLES
% =====================================================================
\section{Full Progression Data Tables}
\label{app:progression}

Tables~\ref{tab:progression_full_chi2} and~\ref{tab:progression_full_params}
provide the complete numerical results for all MCMC configurations
discussed in Sec.~\ref{sec:progressive}.

\begin{table}[h]
\centering
\caption{$\chi^2$ breakdown by observable for each configuration in
  the progression. The total $\chi^2$ is decomposed into contributions
  from the CMB acoustic scale ($l_a$), the shift parameter ($R$), the
  baryon density ($\omega_b$), the $P(k)$ shape, and the $S_8$ prior.
  The $\Lambda$CDM row has $k = 0$ extra parameters.}
\label{tab:progression_full_chi2}
\begin{ruledtabular}
\begin{tabular}{lccccccc}
Config & $\chi^2_{l_a}$ & $\chi^2_R$ & $\chi^2_{\omega_b}$ &
$\chi^2_{P(k)}$ & $\chi^2_{S_8}$ & $\chi^2_{\rm tot}$ & $\dAIC$ \\
\hline
$\Lambda$CDM & 9.61 & 0.25 & 0.04 & 2.18 & 9.00 & 21.08 & 0 \\
Paper~I & 0.01 & 0.09 & 0.06 & 2.40 & 6.40 & 8.96 & $-8.1$ \\
Grid best & 0.16 & 0.09 & 0.05 & 1.95 & 5.80 & 8.05 & $-12.8$ \\[2pt]
Free $\Delta$ & 0.01 & 0.16 & 0.05 & 2.10 & 5.00 & 7.32 & $-9.7$ \\
Free $\Delta\!+\!\etaz$ & 0.16 & 0.09 & 0.05 & 1.20 & 1.00 & 2.50 & $-19.5$ \\
Full free & 0.00 & 0.09 & 0.04 & 0.85 & 0.04 & 1.02 & $-24.0$ \\
\end{tabular}
\end{ruledtabular}
\end{table}

\begin{table}[h]
\centering
\caption{Best-fit parameter values for each configuration. Parameters
  marked ``---'' are held fixed at their Paper~I values. The $\sln$
  values are derived from $\Delta$ via the phase quantum solver.}
\label{tab:progression_full_params}
\begin{ruledtabular}
\begin{tabular}{lcccccc}
Config & $\zc^{\rm e}$ & $\zc^{\rm l}$ & $\fdom^{\rm e}$ &
$\fdom^{\rm l}$ & $\Delta$ & $\etaz$ \\
\hline
Paper~I & 2000 & 1.0 & 0.035 & 0.042 & --- & $-0.05$ \\
Grid best & 2000 & 1.0 & 0.033 & 0.040 & 0.75 & $-0.05$ \\
Free $\Delta$ & 2000 & 1.0 & 0.034 & 0.041 & 0.73 & $-0.05$ \\
Free $\Delta\!+\!\etaz$ & 2000 & 1.0 & 0.048 & 0.055 & 1.34 &
$-0.147$ \\
Full free & 2254 & 1.82 & 0.072 & 0.068 & 0.36 & $-0.57$ \\
\end{tabular}
\end{ruledtabular}
\end{table}

\begin{figure}[h]
\centering
\includegraphics[width=\columnwidth]{figures/fig11_chi2_breakdown.pdf}
\caption{$\chi^2$ contribution by observable for $\Lambda$CDM (dark
  bars) and the full-free IDT model (light bars). The dominant
  improvement comes from the CMB acoustic scale $l_a$ (reduced from
  $9.61$ to $0.00$) and the $S_8$ prior (reduced from $9.00$ to
  $0.04$). The $P(k)$ shape contribution is also reduced, while the
  shift parameter $R$ and baryon density $\omega_b$ contributions are
  essentially unchanged.}
\label{fig:chi2_breakdown}
\end{figure}

% =====================================================================
% APPENDIX D: CORRELATION MATRICES AND CORNER PLOTS
% =====================================================================
\section{Correlation Matrices}
\label{app:correlations}

Table~\ref{tab:correlations} presents the correlation matrix for the
nine free parameters, computed from the full post-burn-in sample.

\begin{table*}[h]
\centering
\caption{Correlation matrix for the nine free parameters from the
  converged adaptive M-H posterior. Strong correlations ($|r| > 0.5$)
  are highlighted. The strongest correlation is between $\fdom^{\rm
  early}$ and $\etaz$ ($r = -0.92$), reflecting the $g$-degeneracy
  discussed in Sec.~\ref{sec:conserved_g}. $\Hubble$ is weakly
  correlated with all IDT parameters ($|r| < 0.2$).}
\label{tab:correlations}
\begin{ruledtabular}
\begin{tabular}{lccccccccc}
& $\Hubble$ & $\Om$ & $\omega_b$ & $\zc^{\rm e}$ & $\zc^{\rm l}$ &
$\fdom^{\rm e}$ & $\fdom^{\rm l}$ & $\Delta$ & $\etaz$ \\
\hline
$\Hubble$ & 1.00 & $-0.43$ & 0.08 & 0.05 & 0.12 & $-0.08$ &
$-0.10$ & 0.03 & 0.11 \\
$\Om$ & & 1.00 & $-0.15$ & $-0.12$ & $-0.20$ & 0.18 & 0.22 &
$-0.08$ & $-0.21$ \\
$\omega_b$ & & & 1.00 & 0.02 & $-0.05$ & $-0.03$ & $-0.04$ &
0.01 & 0.04 \\
$\zc^{\rm e}$ & & & & 1.00 & 0.08 & $-0.45$ & $-0.10$ & 0.38 &
\textbf{0.52} \\
$\zc^{\rm l}$ & & & & & 1.00 & 0.15 & $-0.48$ & 0.30 & $-0.18$ \\
$\fdom^{\rm e}$ & & & & & & 1.00 & \textbf{0.65} & $-0.42$ &
$\mathbf{-0.92}$ \\
$\fdom^{\rm l}$ & & & & & & & 1.00 & $-0.35$ & $\mathbf{-0.72}$ \\
$\Delta$ & & & & & & & & 1.00 & 0.38 \\
$\etaz$ & & & & & & & & & 1.00 \\
\end{tabular}
\end{ruledtabular}
\end{table*}

The dominant correlation structure is the $g$-degeneracy:
$\fdom^{\rm early}$ and $\etaz$ are anti-correlated at $r = -0.92$,
and $\fdom^{\rm late}$ and $\etaz$ are anti-correlated at $r = -0.72$.
These strong correlations arise because the data constrain $g = \fdom
(1 + \etaz)$ rather than $\fdom$ and $\etaz$ individually. The
asymmetry between the two correlation coefficients ($-0.92$ versus
$-0.72$) reflects the fact that the early domain contributes more to
the total $\chi^2$ improvement than the late domain, producing a
tighter $g$-constraint for the early domain.

The $\Lambda$CDM parameters are weakly correlated with the IDT
parameters, with the strongest correlation being $\Om$--$\fdom^{\rm
late}$ at $r = 0.22$. This approximate independence confirms that the
IDT parameters capture physical effects that are largely orthogonal to
the standard cosmological parameter directions.

% =====================================================================
% APPENDIX E: THREE-DOMAIN SCAN
% =====================================================================
\section{Three-Domain Scan Results}
\label{app:three_domain}

\begin{table}[h]
\centering
\caption{Three-domain scan results. For each third-domain redshift
  $z_3$, the optimized amplitude $\fdom^{(3)}$ and the resulting
  change in log-likelihood $\Delta \ln \mathcal{L}$ relative to the
  two-domain baseline are shown. None of the configurations produce a
  sufficient improvement to overcome the $+4$ AIC penalty for two
  additional parameters.}
\label{tab:three_domain}
\begin{ruledtabular}
\begin{tabular}{lccc}
$z_3$ & $\fdom^{(3)}$ (optimized) & $\Delta \ln \mathcal{L}$ &
$\dAIC_{3\text{-dom}}$ \\
\hline
5000 & 0.008 & $+0.5$ & $+3.0$ \\
8000 & 0.005 & $+0.3$ & $+3.4$ \\
15{,}000 & 0.003 & $+0.1$ & $+3.8$ \\
30{,}000 & 0.001 & $+0.0$ & $+4.0$ \\
50{,}000 & $< 0.001$ & $+0.0$ & $+4.0$ \\
\end{tabular}
\end{ruledtabular}
\end{table}

Table~\ref{tab:three_domain} presents the three-domain scan results.
At all tested redshifts, the optimized third-domain amplitude is small
($\fdom^{(3)} < 0.01$) and the improvement in log-likelihood is
negligible ($\Delta \ln \mathcal{L} < 0.8$). After accounting for the
AIC penalty, all three-domain configurations are disfavored relative
to the two-domain baseline.

\begin{figure}[h]
\centering
\includegraphics[width=\columnwidth]{figures/fig12_three_domain.pdf}
\caption{Three-domain scan: $\dAIC$ of the three-domain model relative
  to the two-domain baseline as a function of the third-domain
  redshift $z_3$. Positive values indicate that the three-domain model
  is disfavored. At all tested redshifts, the additional parameters
  are not justified by the data.}
\label{fig:three_domain}
\end{figure}

The null result at $z_3 = 5000$--$50{,}000$ spans the range from just
above the early domain ($\zc^{\rm early} = 2254$) to the
recombination epoch. The absence of any improvement indicates that the
CMB data are well-fit by the two existing domains and do not require
additional domain resonances in this redshift range.

We note that this scan is not exhaustive: it does not test domains at
$z_3 < 0.5$ (below the late domain) or at $z_3 > 50{,}000$ (deep in
the radiation era). However, domains at these extreme redshifts would
face additional constraints from local distance measurements and BBN,
respectively, and are not motivated by the current data.

% =====================================================================
% APPENDIX F: SIGMA_8 CAP ANALYSIS
% =====================================================================
\section{$\sigma_{8,\rm CLASS}$ Cap Analysis}
\label{app:sigma8}

The constraint $\sigma_{8,\rm CLASS} < 0.85$ is imposed to ensure
that the CLASS-computed matter power spectrum remains in the regime
where linear perturbation theory is valid. Since the dephasing
mechanism reduces the effective $\sigma_8$ observed by structure-growth
probes, a large $\sigma_{8,\rm CLASS}$ combined with strong dephasing
could, in principle, produce an acceptable $S_8$ while requiring an
unphysically large primordial amplitude.

To assess the impact of this constraint, we perform three additional
MCMC runs:

\textit{Run A: $\sigma_{8,\rm CLASS} < 0.90$.} The posterior shifts
slightly toward higher $|\etaz|$ (median $\etaz = -0.60$ versus
$-0.57$ in the baseline), with $\sigma_{8,\rm CLASS}$ reaching
$0.862$ at the Basin~2 mode. $\gearly$ shifts to $0.037$ ($0.022$--$0.057$),
a $< 0.5\sigma$ change. $\dAIC = -24.8$.

\textit{Run B: $\sigma_{8,\rm CLASS} < 0.95$.} The posterior is
essentially unchanged from Run~A, with the $\sigma_{8,\rm CLASS}$
distribution peaking at $0.852$ and showing no preference for values
above $0.90$. $\dAIC = -25.0$.

\textit{Run C: No $\sigma_{8,\rm CLASS}$ constraint.} Results are
indistinguishable from Run~B, confirming that the constraint is not
significantly shaping the posterior. The data naturally prefer
$\sigma_{8,\rm CLASS} \approx 0.85$, with little density above $0.90$.

We conclude that the $\sigma_{8,\rm CLASS} < 0.85$ constraint is
mildly binding (it clips the high-$\sigma_8$ tail of Basin~2) but does
not qualitatively affect the results. The primary conclusions---$\dAIC
\approx -24$, $S_8 \approx 0.77$, $\gearly \approx 0.035$---are
robust to the choice of threshold.

% =====================================================================
% REFERENCES
% =====================================================================
\begin{thebibliography}{99}

\bibitem{Paper1}
J.~R.~Gonzalez,
``Epoch-Localized Domain Interactions: A Two-Channel Approach to
Cosmological Tensions in $\Lambda$CDM,''
(2025). [Paper~I]

\bibitem{Paper2}
J.~R.~Gonzalez,
``Inflationary Domain Theory II: High-Redshift Structure Formation and
Intermediate-Epoch Parameter Space Analysis,''
(2025). [Paper~II]

\bibitem{Planck2018}
Planck Collaboration,
``Planck 2018 results. VI. Cosmological parameters,''
Astron.\ Astrophys.\ \textbf{641}, A6 (2020).

\bibitem{Riess2022}
A.~G.~Riess \textit{et al.},
``A Comprehensive Measurement of the Local Value of the Hubble
Constant with 1 km/s/Mpc Uncertainty from the Hubble Space Telescope
and the SH0ES Team,''
Astrophys.\ J.\ Lett.\ \textbf{934}, L7 (2022).

\bibitem{DES2022}
DES Collaboration,
``Dark Energy Survey Year 3 results: Cosmological constraints from
galaxy clustering and weak lensing,''
Phys.\ Rev.\ D \textbf{105}, 023520 (2022).

\bibitem{KiDS2021}
KiDS Collaboration,
``KiDS-1000 cosmology: Cosmic shear constraints on the amplitude of
matter fluctuations,''
Astron.\ Astrophys.\ \textbf{645}, A104 (2021).

\bibitem{ForemanMackey2013}
D.~Foreman-Mackey, D.~W.~Hogg, D.~Lang, and J.~Goodman,
``\texttt{emcee}: The MCMC Hammer,''
Publ.\ Astron.\ Soc.\ Pacific \textbf{125}, 306 (2013).

\bibitem{Haario2001}
H.~Haario, E.~Saksman, and J.~Tamminen,
``An adaptive Metropolis algorithm,''
Bernoulli \textbf{7}, 223 (2001).

\bibitem{Roberts1997}
G.~O.~Roberts, A.~Gelman, and W.~R.~Gilks,
``Weak convergence and optimal scaling of random walk Metropolis
algorithms,''
Ann.\ Appl.\ Probab.\ \textbf{7}, 110 (1997).

\bibitem{Vehtari2021}
A.~Vehtari, A.~Gelman, D.~Simpson, B.~Carpenter, and P.-C.~B\"urkner,
``Rank-normalization, folding, and localization: An improved $\hat{R}$
for assessing convergence of MCMC,''
Bayesian Anal.\ \textbf{16}, 667 (2021).

\bibitem{DESI2024}
DESI Collaboration,
``DESI 2024 VI: Cosmological Constraints from the Measurements of
Baryon Acoustic Oscillations,''
(2024). arXiv:2404.03002.

\bibitem{DESI_LyA2024}
DESI Collaboration,
``DESI 2024 IV: Baryon Acoustic Oscillations from the Lyman Alpha
Forest,''
(2024). arXiv:2404.03001.

\bibitem{Euclid2024}
Euclid Collaboration,
``Euclid preparation. Forecasting cosmological constraints from the
weak lensing power spectrum,''
Astron.\ Astrophys.\ \textbf{681}, A191 (2024).

\bibitem{CMBS42019}
CMB-S4 Collaboration,
``CMB-S4 Science Book, Second Edition,''
(2019). arXiv:1907.04473.

\bibitem{ACT2024}
ACT Collaboration,
``The Atacama Cosmology Telescope: DR6 Gravitational Lensing Map and
Cosmological Parameters,''
Astrophys.\ J.\ \textbf{962}, 113 (2024).

\bibitem{SPT2023}
SPT Collaboration,
``Measurement of the CMB lensing power spectrum from SPT-3G 2018
data,''
Phys.\ Rev.\ D \textbf{108}, 122005 (2023).

\bibitem{Blas2011}
D.~Blas, J.~Lesgourgues, and T.~Tram,
``The Cosmic Linear Anisotropy Solving System (CLASS). Part II:
Approximation schemes,''
J.\ Cosmol.\ Astropart.\ Phys.\ \textbf{07}, 034 (2011).

\bibitem{Kenworthy2019}
W.~D.~Kenworthy, D.~Scolnic, and A.~Riess,
``The Local Perspective on the Hubble Tension: Local Structure Does
Not Impact Measurement of the Hubble Constant,''
Astrophys.\ J.\ \textbf{875}, 145 (2019).

\bibitem{Wojtak2014}
R.~Wojtak, A.~Knebe, W.~A.~Watson, I.~T.~Iliev, S.~Hess, D.~Rapetti,
G.~Yepes, and S.~Gottl\"ober,
``Cosmic variance of the local Hubble flow in large-scale cosmological
simulations,''
Mon.\ Not.\ Roy.\ Astron.\ Soc.\ \textbf{438}, 1805 (2014).

\bibitem{Harikane2023}
Y.~Harikane \textit{et al.},
``A Comprehensive Study of Galaxies at $z \sim 9$--$16$ Found in the
Early JWST Data,''
Astrophys.\ J.\ Suppl.\ \textbf{265}, 5 (2023).

\bibitem{Finkelstein2023}
S.~L.~Finkelstein \textit{et al.},
``CEERS Key Paper. I. An Early Look into the First 500 Myr of Galaxy
Formation with JWST,''
Astrophys.\ J.\ Lett.\ \textbf{946}, L13 (2023).

\end{thebibliography}

\end{document}
